% Part: sets-functions-relations
% Chapter: functions
% Section: kinds

\documentclass[../../../include/open-logic-section]{subfiles}

\begin{document}

\olfileid{sfr}{fun}{kin}
\olsection{函数的种类}

\begin{explain}
为我们最常遇到的几类函数引入一种分类法会很有用。

首先，我们可以考察具有如下性质的函数：其陪域中的每个元素都是该函数的值。这样的函数称为满射的，如\olref{fig:surjective}所示。

\begin{figure}
  \olasset{assets/diagrams/surjective.tikz}
  \caption{满射函数的陪域中的每个元素都是该函数的值。}
  \ollabel{fig:surjective}
\end{figure}
\end{explain}

\begin{defn}[满射函数]
函数 $f \colon A \rightarrow B$ 是\emph{满射的}，当且仅当 $B$ 也是~$f$ 的值域，即对每个 $y \in B$，都至少有一个 $x \in A$ 使得~$f(x) = y$，用符号表示：
\[
  (\forall y \in B)(\exists x \in A)f(x) = y.
\]
我们称这样的函数为从 $A$ 到 $B$ 的满射。
\end{defn}

\begin{explain}
如果要证明 $f$ 是满射，你需要证明 $f$ 的陪域中的每个对象都是某个输入 $x$ 对应的值 $f(x)$。

注意，任何函数都\emph{诱导}出一个满射。事实上，给定函数 $f \colon A \to B$，令 $f' \colon A \to \ran{f}$ 定义为 $f'(x) = f(x)$。由于 $\ran{f}$ 被\emph{定义}为 $\Setabs{f(x) \in B}{x \in A}$，该函数 $f'$ 必定是一个满射。
\end{explain}

\begin{explain}
任何函数都将每个可能的输入映射到唯一的输出。但也存在一些函数，它们从不会把不同的输入映射到相同的输出。这样的函数称为单射的，如\olref{fig:injective}所示。
\begin{figure}
  \olasset{assets/diagrams/injective.tikz}
  \caption{单射函数从不会把两个不同的自变元映射到同一个值。}
  \ollabel{fig:injective}
\end{figure}
\end{explain}

\begin{defn}[单射函数]
函数 $f \colon A \rightarrow B$ 是\emph{单射的}，当且仅当对每个 $y \in B$，都至多有一个 $x \in A$ 使得~$f(x) = y$。我们称这样的函数为从 $A$ 到~$B$ 的单射。
\end{defn}

\begin{explain}
如果要证明 $f$ 是单射，你需要证明：对 $f$ 定义域中的任意元素 $x$ 和 $y$，如果 $f(x)=f(y)$，则 $x=y$。
\end{explain}

\begin{ex}
由 $f(x) = 1$ 给出的常值函数 $f\colon \Nat \to \Nat$ 既非单射也非满射。

由 $f(x) = x$ 给出的恒等函数 $f\colon \Nat \to \Nat$ 既是单射的又是满射的。

由 $f(x) = x+1$ 给出的后继函数 $f \colon \Nat \to \Nat$ 是单射的，但不是满射的。

函数 $f \colon \Nat \to \Nat$ 定义如下：
\[
  f(x) =
  \begin{cases}
    \frac{x}{2} & \text{若 $x$ 是偶数} \\
    \frac{x+1}{2} & \text{若 $x$ 是奇数。}
  \end{cases}
\]
该函数是满射的，但不是单射的。
\end{ex}

\begin{explain}
很多时候，我们要考虑既是单射又是满射的函数。我们称这样的函数为双射的。它们类似于\olref{fig:bijective}中所描绘的函数。双射有时也被称为\emph{一一对应}，因为它们将陪域中的元素与定义域中的元素唯一地配对。
\begin{figure}
  \olasset{assets/diagrams/bijective.tikz}
  \caption{双射函数将陪域中的元素与定义域中的元素唯一地配对。}
  \ollabel{fig:bijective}
\end{figure}
\end{explain}

\begin{defn}[双射]
函数 $f \colon A \to B$ 是\emph{双射的}，当且仅当它既是满射的又是单射的。我们称这样的函数为从 $A$ 到~$B$（或在 $A$ 与~$B$ 之间）的双射。
\end{defn}

\end{document}